\documentclass[11pt,a4paper]{article}
\usepackage[T1]{fontenc}
\usepackage[utf8]{inputenc}
\usepackage{lmodern,microtype}
\usepackage[margin=26mm]{geometry}
\usepackage{amsmath,amssymb,amsthm,mathtools}
\usepackage{booktabs,longtable,array,enumitem}
\usepackage{xcolor,fancyhdr}
\usepackage[hidelinks]{hyperref}
\hypersetup{pdftitle={Exact Hankel Periods, Parking-Outcome Formulas, and Three Further Investigations},pdfauthor={KEIO}}
\pagestyle{fancy}
\fancyhf{}
\fancyhead[L]{\small T96--T100: proofs and research boundaries}
\fancyhead[R]{\small KEIO}
\fancyfoot[C]{\thepage}
\setlength{\headheight}{14pt}
\setlength{\emergencystretch}{2em}
\allowdisplaybreaks
\newtheorem{theorem}{Theorem}[section]
\newtheorem{lemma}[theorem]{Lemma}
\newtheorem{proposition}[theorem]{Proposition}
\newtheorem{corollary}[theorem]{Corollary}
\theoremstyle{definition}
\newtheorem{definition}[theorem]{Definition}
\theoremstyle{remark}
\newtheorem{remark}[theorem]{Remark}
\newcommand{\F}{\mathbb F}
\newcommand{\Q}{\mathbb Q}
\newcommand{\Z}{\mathbb Z}
\newcommand{\C}{\mathbb C}
\DeclareMathOperator{\Aut}{Aut}
\DeclareMathOperator{\Cay}{Cay}
\begin{document}
\title{Exact Hankel Periods, Parking-Outcome Formulas,\\
and Three Further Investigations}
\author{KEIO}
\date{T96--T100\quad\textbar\quad27 October 2026}
\maketitle
\begin{abstract}
We prove Han's least-period conjecture for every shifted binary
Rudin--Shapiro Hankel sequence with shift at least three. The proof
combines explicit norm-one units in a quadratic function field with an
approximation-error criterion: a sufficiently long gap between normal
indices forces a square norm, and the possible finite pole orders then
force that norm to be one. This identifies the exact spacing of long
gaps and proves minimality, in addition to periodicity.
We also give nonrecursive coefficient-extraction formulas answering two
parking-outcome enumeration questions of Ferreri, Harris, Martinez, and
Swartz, including arbitrary repeated capacities and every position of
a single unavailable spot. Three further investigations concern Stern
parity counts, rank-two average mixing, and a Legendre-determinant
nonvanishing conjecture. Their incomplete conclusions are explicitly
separated from the complete original-target proofs. Written arguments,
primary-source comparisons, and exact reproducibility programs accompany
the results. Worldwide priority and a historically unprecedented
foundational principle are not certified.
\end{abstract}

\section*{Scope and success criteria}
The five tracks below are research attempts, not five asserted solutions
of open problems. The original targets are identified before the
auxiliary constructions, and their parameter ranges and mathematical
objects are retained. A complete solution, a substantive new argument,
and historical priority are separate assertions.

The main complete proof is T98: Han's 2016 Conjecture~4.2, including
the least-period requirement for every shift. T97 answers two explicit
questions in Section~6 of the 2025 parking-outcome paper; those questions
have author-and-section provenance but no individual conjecture numbers.
The terminal-event and increasing-run decompositions used there extend
classical methods. The three other tracks must be read with their stated
remaining gaps; their auxiliary statements do not replace their original
targets.

The requested combination of an unchanged previously open named problem,
an unconditional complete proof, a historically unprecedented principle,
and established absence of an identical prior resolution is not certified
for all five tracks, or asserted as an absolute historical certification
for any one track. Bounded literature searches are reported as such.
Internal mathematical reviews and exact calculations support the proofs;
they are not external peer review or proof-assistant formalizations.

This round follows the preceding T91--T95 investigation. Its previously
proved AFM lucky-set identity and zero-residue Stern growth theorem are
not counted again. The supplied mathematical corpus is treated as a
collection of research drafts, consistent with its own provenance
statement; original primary papers are used to establish the external
targets and the published inputs.


\clearpage
\section{T96: a degree-preserving ternary fold for odd Stern coefficients}
\label{sec:t96}

The unchanged target is Question~2.10 of Ulas
\cite{t96:ulas}: for $r\in\{0,1\}$, does
$\Pi_{r,2}(x)/(\log_2x)^2$ tend to infinity? Here
\[
 B_0=0,\quad B_1=1,\quad B_{2n}=tB_n,\quad
 B_{2n+1}=B_n+B_{n+1},\qquad e(n)=\deg B_n,
\]
and
\[
 \Pi_{r,2}(x)=\#\left\{1\le n\le x:
 B_n(t)\equiv1+r\sum_{i=1}^{e(n)}t^i\pmod2\right\}.
\]
The degree in this definition is the degree over $\mathbb Z$, not the
possibly smaller degree after reduction modulo two. The earlier T92
argument treated $r=0$. The construction below advances $r=1$, but its
exponent is less than two, so it does not close the original question.

\begin{theorem}[A lower bound in the residue-one branch]
Put $\alpha=\log_2 3$. For $L=\lfloor\log_2x\rfloor\ge3$,
\[
 \Pi_{1,2}(x)\ge
 3^{\left\lfloor\log_2((L+2)/5)\right\rfloor}
 \ge \frac13\left(\frac{L+2}{5}\right)^\alpha.
\]
In particular $\Pi_{1,2}(x)\gg(\log x)^{\log_2 3}$, and
$\Pi_{1,2}(x)/(\log x)^A\to\infty$ for every fixed
$A<\log_2 3$.
\end{theorem}

The auxiliary problem is to find several reversible operations on binary
words which preserve the full integer-degree congruence, rather than
only its reduction. A parity identity gives three candidate operations;
an integer leading-coefficient calculation proves that none creates
hidden even coefficients in higher degrees.

For a binary word $w$, write $[w]_2$ for its value and
$\operatorname{rev}(w)$ for its reversal. Products of the digit matrices
\[
 L=\begin{pmatrix}t&0\\1&1\end{pmatrix},\qquad
 R=\begin{pmatrix}1&1\\0&t\end{pmatrix}
\]
act on $v=(0,1)^{\mathsf T}$ in the order in which digits are read.
Thus $M_{ab}=M_bM_a$ and
$M_wv=(B_{[w]_2},B_{[w]_2+1})^{\mathsf T}$.

\begin{lemma}[The ternary fold]
Suppose that $w$ has at least three digits, starts and ends in $11$, and
\[
 B_{[w]_2}(t)\equiv\sum_{i=0}^{E}t^i\pmod2,
 \qquad E=e([w]_2).
\]
Remove its last digit, writing $w=g1$. For $k=0,1,2$, define
\[
 \Phi_k(w)=g\,\operatorname{rev}(g)\,0^k11.
\]
Each resulting word starts and ends in $11$, and
\[
 e([\Phi_k(w)]_2)=2E+1,\qquad
 B_{[\Phi_k(w)]_2}(t)\equiv\sum_{i=0}^{2E+1}t^i\pmod2.
\]
The three maps have disjoint images, and each has an explicit inverse
on its image.
\end{lemma}

\begin{proof}
First work over $\mathbb F_2[t]$. Put $d=t+1$ and
$K=\left(\begin{smallmatrix}d&1\\1&d\end{smallmatrix}\right)$.
Direct multiplication gives $KL=L^{\mathsf T}K$ and
$KR=R^{\mathsf T}K$. Hence
\[
 KM_{\operatorname{rev}(g)}=M_g^{\mathsf T}K.
\]
Also $(1,0)R^2=(1,d)=v^{\mathsf T}K$ and
$(1,d)L^k=(1,d)$ for every $k\ge0$. If $M_gv=(a,b)^{\mathsf T}$,
then $B_{[w]_2}=a+b$, and therefore
\begin{align*}
 B_{[\Phi_k(w)]_2}
 &=v^{\mathsf T}KM_{\operatorname{rev}(g)}M_gv\\
 &=(a,b)K(a,b)^{\mathsf T}
   =d(a+b)^2
   =\sum_{i=0}^{2E+1}t^i.
\end{align*}
The last equality uses
$(1+t)(1+t+\cdots+t^E)^2=1+t+\cdots+t^{2E+1}$ in
characteristic two. This parity calculation alone would not suffice.

We now bound the degrees before reduction. All digit-matrix entries
have nonnegative integer coefficients. The degree recurrence implies
$|e(n+1)-e(n)|\le1$ for $n\ge1$, by induction. It also implies
$e(n+1)\ge e(n)$ when $n$ is odd. Since $g$ ends in $1$,
\[
 E=e([g1]_2)=e([g]_2+1).
\]
Write $M_g=\left(\begin{smallmatrix}p&q\\r&s\end{smallmatrix}\right)$.
Thus $\deg q\le E$ and $\deg s=E$. Because $g$ begins with $11$,
we can write $M_g=A R^2$, where $A$ again has nonnegative coefficients.
The second column of $AR^2$ equals $d$ times its first column plus
$t^2$ times the second column of $A$. Consequently
\[
 \deg p\le E-1,\qquad \deg r\le E-1.
\]
We use the convention $\deg0=-\infty$ here.

Over $\mathbb Z[t]$ the exact self-adjointness matrix is
\[
 K_{\mathbb Z}=\begin{pmatrix}1-t&1\\1&1-t\end{pmatrix},
 \qquad
 K_{\mathbb Z}M_{\operatorname{rev}(g)}=M_g^{\mathsf T}K_{\mathbb Z}.
\]
Its leading matrix is $-I$. Comparing the highest powers of $t$ on
the two sides shows that $M_{\operatorname{rev}(g)}$ has the same
maximum degree $E$ as $M_g$, with transposed leading-coefficient
matrix. The first column of $M_g$ has no term of degree $E$;
the first row of $M_{\operatorname{rev}(g)}$ therefore has degree
at most $E-1$. It follows that, on putting
\[
 M_{\operatorname{rev}(g)}M_gv=(A_0,B_0)^{\mathsf T},
\]
one has $\deg A_0\le2E-1$ and $\deg B_0\le2E$.

Let $S_k=1+t+\cdots+t^{k-1}$, with $S_0=0$. Over the integers,
\[
 (1,0)R^2L^k=(t^k+dS_k,d).
\]
For $0\le k\le2$ the first entry has degree at most two. Thus
\[
 \deg B_{[\Phi_k(w)]_2}
 \le\max(2E-1+k,2E+1)=2E+1.
\]
The parity calculation already gives a nonzero coefficient in degree
$2E+1$, proving equality and the full required congruence.

Finally, $g\operatorname{rev}(g)$ ends in $1$. For $k=0$ the output
ends in at least three ones; for $k=1$ or $2$, the final run $11$ is
preceded by exactly $k$ zeros. These distinguish the three images.
After deleting $0^k11$, take the first half of the remaining word to
recover $g$, and append $1$ to recover $w$. This proves injectivity.
\end{proof}

\begin{proof}[Proof of the theorem]
Start with $w_0=111$, so $B_7=1+t+t^2$. For every string
$(k_1,\ldots,k_j)\in\{0,1,2\}^j$, apply the corresponding folds.
Their inverses give exactly $3^j$ distinct binary words at level $j$.
All of them satisfy the original residue-one congruence, with
\[
 e([w_j]_2)=3\cdot2^j-1.
\]
The lengths satisfy $\ell_0=3$ and
$\ell_i=2\ell_{i-1}+k_i$. In particular,
\[
 \ell_j=3\cdot2^j+\sum_{i=1}^j2^{j-i}k_i
 \le5\cdot2^j-2.
\]
For $j=\lfloor\log_2((L+2)/5)\rfloor$, all these words have at
most $L$ digits, and hence represent indices below $2^L\le x$.
This proves the first bound. The second uses
$3^{\lfloor y\rfloor}\ge3^{y-1}$.
\end{proof}

\paragraph{What is and is not resolved.}
The source proves the lower bound $\Pi_{1,2}(x)\ge\log_2x$ in
Corollary~2.7. The exponent $\log_2 3$ above is a stronger quantitative
bound for exactly the same counting function, with its integer-degree
condition intact. It is still smaller than the exponent two in
Question~2.10, so this is a partial advance, not a complete answer.
The preceding round's $r=0$ result is not counted again as a new result.

Digit matrices, reversal symmetries, and characteristic-two quadratic
identities are established tools. Related reversal work includes
Spiegelhofer~\cite{t96:spiegelhofer}, for a different polynomial
generalization of Stern's sequence. The specific contribution proved
here is the three-branch operation together with the integer leading-
degree argument needed for Ulas's exact condition. The bounded search
did not locate this construction or its stated lower bound; that is not
a certification of a historically unprecedented principle or worldwide
priority. This track does not meet the requested full-success gate.

The program \path{t96_verify.py} checks the integer self-adjointness
identities, the entire ternary construction through level six, additional
original residue-one seeds, and independent integer polynomial examples.
It includes negative controls showing why extending $k$ to three or
dropping the initial-$11$ restriction can introduce a degree gap.

\clearpage
\section{T97: terminal-event enumeration and two parking-outcome formulas}
\label{sec:t97parking}

The targets are the first two enumeration questions in Section~6 of
Ferreri--Harris--Martinez--Swartz \cite{t97fhms}. The first asks for the
number of outcomes admitting $k$ lucky cars for an arbitrary capacity
vector, including repetitions. The second asks for a closed formula for
$c_s(n,I)$, the number of outcomes with a fixed lucky set $I$ when the
single unavailable spot is $s$. These questions are explicit in the
pre-existing source, although they are not assigned conjecture numbers.
We give nonrecursive coefficient-extraction formulas over their full
parameter ranges. The outcomes, rather than the preference lists, are
counted, and an outcome admitted by several lucky sets is counted once.

The auxiliary device is to order the terminal occurrences of letters
and enumerate the intervals between them. It yields a rational series
with only $r!$ summands for $r$ available spots, independent of their
capacities. A separate increasing-run decomposition gives all values of
$c_s(n,I)$ simultaneously from $2^{|I|}$ products. These are complete
enumerative formulas, but the use of classical order decompositions and
formal geometric series does not establish the requested historically
unprecedented foundational principle.

\subsection{An auxiliary terminal-event problem}

Fix $r$ letters, each required to occur at least once in a finite word.
The final occurrence of letter $j$ is its terminal event. Fix their order
$\sigma=(\sigma(1),\ldots,\sigma(r))$, and write $\rho(j)$ for the rank
of $j$ in this order. The $h$th interval, for $0\le h<r$, consists of the
ordinary occurrences after the first $h$ terminal events and before the
next terminal event. Only letters with $\rho(j)>h$ can occur in this
interval. Give an ordinary occurrence of $j$ in interval $h$ weight
$x_j w_{j,h}$, and its terminal occurrence weight $a_j$.

\begin{lemma}[Terminal-event generating series]
For a fixed terminal order, the generating series of these words is
\[
 \frac{\prod_{j=1}^r a_j}
 {\displaystyle\prod_{h=0}^{r-1}
       \left(1-\sum_{\rho(j)>h}x_jw_{j,h}\right)}.
\]
Consequently the coefficient of $\prod_jx_j^{m_j-1}$ counts words
with $m_j$ occurrences of each letter, with the stated weights.
\end{lemma}

\begin{proof}
The word has a unique decomposition
\[
 W_0\,\sigma(1)\,W_1\,\sigma(2)\cdots
 W_{r-1}\,\sigma(r),
\]
where $W_h$ is any word on the letters whose terminal events have not
yet occurred. The ordinary generating series of $W_h$ is the geometric
series of the sum of its allowed letter weights. Multiplying these
series and the terminal weights proves the formula. Conversely every
such choice has precisely the indicated terminal order: its terminal
occurrence is present and the letter is forbidden afterward. Thus the
decomposition is bijective, including empty intervals.
\end{proof}

\subsection{Arbitrary capacities: the full outcome-counting formula}

Write the distinct available spots as $v_1<\cdots<v_r$, with positive
capacities $m_1,\ldots,m_r$ and $n=\sum_jm_j$. Thus the vector
$\boldsymbol u$ contains $m_j$ copies of $v_j$. Every outcome is an
ordered set partition $B=(B_1,\ldots,B_r)$ of $[n]$ with $|B_j|=m_j$.
All such partitions are possible outcomes, by letting each car prefer
its designated spot. Define
\[
 F(B)=\bigcup_{v_j=1} B_j\ \cup
 \bigcup_{\substack{j\ge2\\v_j=v_{j-1}+1}}
       \{i\in B_j:i<\max B_{j-1}\}.
\]
This is the set of cars forced to be lucky by the outcome. The following
criterion is the content of the source's Theorem~3.7; we include the
short verification to specify precisely what is counted.

\begin{lemma}[The possible lucky sets of one outcome]
An outcome $B$ admits the exact lucky set $I$ if and only if
$F(B)\subseteq I\subseteq[n]$.
\end{lemma}

\begin{proof}
A car parking at spot $1$ must prefer $1$. If a car $i$ parks at $v_j$
while the preceding spot $v_j-1=v_{j-1}$ still has an unoccupied place,
then every preference below $v_j$ stops before $v_j$. Such a car must
be lucky. The preceding spot is not full when $i$ arrives exactly when
$i<\max B_{j-1}$. This proves necessity.

For sufficiency, let each car in $I$ prefer its designated spot $v_j$,
and each car outside $I$ prefer $v_j-1$. The latter is positive: cars at
spot $1$ are in $F(B)$. The preceding spot is either unavailable or is
already full, because its designated cars all have smaller labels.
Induction on the arriving label shows that every car parks in its
designated block and that the exact lucky set is $I$.
\end{proof}

Put
\[
 P_{\boldsymbol u}(q)=\sum_{B}q^{|F(B)|}.
\]
This records the minimum possible number of lucky cars for each outcome,
not the number for a single chosen preference list.
For $\sigma\in\mathfrak S_r$ let $\rho_\sigma(j)$ be the rank of $j$
in $\sigma$, and for $0\le h\le r$ set
\[
 \epsilon_\sigma(j,h)=
 \begin{cases}
 1,&v_j=1,\\
 1,&j\ge2,\ v_j=v_{j-1}+1,\ \rho_\sigma(j-1)>h,\\
 0,&\text{otherwise}.
 \end{cases}
\]
Define the integer $A_\sigma$ and linear forms $Q_{\sigma,h}$ by
\[
 A_\sigma=\sum_{j=1}^r
      \epsilon_\sigma\bigl(j,\rho_\sigma(j)\bigr),\qquad
 Q_{\sigma,h}(\mathbf x,q)
       =\sum_{\rho_\sigma(j)>h}x_jq^{\epsilon_\sigma(j,h)}
       \quad(0\le h<r).
\]

\begin{theorem}[General-capacity outcome formula]
For every weakly increasing positive vector $\boldsymbol u$, with
distinct spots and capacities as above,
\begin{equation}\label{eq:t97general}
 P_{\boldsymbol u}(q)
 = [x_1^{m_1-1}\cdots x_r^{m_r-1}]
     \sum_{\sigma\in\mathfrak S_r}
       \frac{q^{A_\sigma}}
       {\prod_{h=0}^{r-1}(1-Q_{\sigma,h}(\mathbf x,q))}.
\end{equation}
The number of distinct outcomes admitting exactly $k$ lucky cars is
\begin{equation}\label{eq:t97union}
 \left|\bigcup_{\substack{I\subseteq[n]\\|I|=k}}
                \mathcal O_{\boldsymbol u}(I)\right|
       =[q^k]\frac{P_{\boldsymbol u}(q)}{1-q}
       \qquad(0\le k\le n).
\end{equation}
All series in \eqref{eq:t97general} are expanded formally at
$\mathbf x=0$.
\end{theorem}

\begin{proof}
Encode an ordered partition $B$ by the word $w_1\cdots w_n$ where
$w_i=j$ precisely when car $i$ belongs to $B_j$. The terminal event of
letter $j$ occurs at the label $\max B_j$. Its relative order is a
unique permutation $\sigma$.

Suppose an ordinary occurrence of $j$ lies after $h$ terminal events.
It is forced lucky if $v_j=1$, or if its preceding available spot is
adjacent and that spot's maximum label still lies in the future. The
latter condition is exactly $\rho_\sigma(j-1)>h$. Thus its weight is
$x_jq^{\epsilon_\sigma(j,h)}$. At the terminal occurrence of $j$,
the corresponding criterion is
$\rho_\sigma(j-1)>\rho_\sigma(j)$; its weight is therefore
$q^{\epsilon_\sigma(j,\rho_\sigma(j))}$. The auxiliary lemma now
gives the summand for $\sigma$. Extracting the indicated capacities
and summing the disjoint terminal orders proves
\eqref{eq:t97general}.

An outcome with $|F(B)|=a$ admits a lucky set of size $k$ precisely when
$a\le k\le n$, by the preceding lemma. Hence its contribution to the
union count is one for every $k\ge a$. Taking the indicated coefficient
of $(1-q)^{-1}$ proves \eqref{eq:t97union}. In particular, we do not
sum the cardinalities of the overlapping sets
$\mathcal O_{\boldsymbol u}(I)$.
\end{proof}

Formula \eqref{eq:t97general} is a finite, nonrecursive rational
coefficient formula in $r$ variables with $r!$ explicitly specified
terms. It contains no sum over parking functions, outcomes, or choices
of their $n$ labels. For a fixed number $r$ of available spots,
truncation at the requested powers evaluates it using polynomially
many arithmetic operations in $n$; the capacities may be arbitrarily
large. No polynomial-time bound in variable $r$ is asserted. When all capacities
are one, the terminal orders are the outcome permutations themselves;
that already-solved boundary case is not presented as an additional advance.

There is also a useful factorization. Split the available spots into
maximal consecutive clusters, containing $N_1,\ldots,N_b$ cars. Let
$P_\ell(q)$ be the forced-lucky polynomial within the $\ell$th cluster,
with its first spot forced only if that cluster contains spot $1$.
Then
\[
 P_{\boldsymbol u}(q)=\binom{n}{N_1,\ldots,N_b}
                         \prod_{\ell=1}^b P_\ell(q).
\]
Indeed, choose the labels assigned to the clusters, then standardize
their relative orders. The forced-lucky comparisons never cross an
unavailable spot. This reduces the terminal-order sum to a separate sum
for each consecutive cluster.

For example, for spots $1,2$ with capacities $a,b$, a direct
specialization gives
\[
 P_{(1^a,2^b)}(q)
   =q^a\sum_{j=0}^b\binom{a+j-1}{a-1}q^j.
\]
Indeed, if $\max B_1=a+j$, exactly $j$ cars of $B_2$ precede it,
and the other $a-1$ cars of $B_1$ can be chosen in the displayed
number of ways. For capacities $(3,2,2)$ at spots $(1,2,3)$ the formula
gives
\[
 P_{\boldsymbol u}(q)=q^3+8q^4+61q^5+95q^6+45q^7.
\]
Thus the numbers of outcomes admitting $k=3,4,5,6,7$ lucky cars are
$1,9,70,165,210$, respectively, with no duplication between lucky sets.

\subsection{One unavailable spot and an arbitrary fixed lucky set}

Now all spots in $[n+1]\setminus\{s\}$ have capacity one, and
$I\subseteq[n]$ is fixed, with $k=|I|$. As in the source's
Definition~3.14, let $c_s(n,I)$ count their possible outcomes with
exact lucky set $I$. For $J\subseteq I$ and $j\notin I$ put
\[
 \lambda_J(j)=|J\cap[j-1]|,\qquad
 \mu_J(j)=1+|(I\setminus J)\cap[j-1]|.
\]

\begin{theorem}[Closed formula for every unavailable spot]
For every $n\ge1$, $1\le s\le n+1$, and $I\subseteq[n]$,
\begin{equation}\label{eq:t97cut}
 c_s(n,I)=[z^{s-1}]
   \sum_{J\subseteq I}|J|!\,(k-|J|)!\,z^{|J|}
       \prod_{j\in[n]\setminus I}
                  \bigl(\lambda_J(j)z+\mu_J(j)\bigr).
\end{equation}
Thus one polynomial with $2^k$ product summands provides all $n+1$
values, without recursion in $s$ or enumeration of $n!$ outcomes.
\end{theorem}

\begin{proof}
The unavailable spot splits an outcome permutation into a left word
of length $s-1$ and a right word of length $n-s+1$. Inside either word,
every descent bottom must be in $I$. The first left letter, if present,
must be in $I$; the first right letter need not be, because its preceding
spot is unavailable. These conditions are exactly the forced-lucky
criterion already proved.

Let $J$ be the lucky labels on the left. Split the left word immediately
before each label in $J$. Each resulting run starts with that lucky
label and then contains only unlucky labels in increasing order, all
larger than its first label. The $|J|$ runs have an arbitrary order.
On the right, split immediately before each lucky label in
$I\setminus J$. There is also a possibly empty initial run of unlucky
labels. Regard this initial run as having a fictitious first label $0$;
it stays first, while the other $k-|J|$ runs have an arbitrary order.

For each unlucky label $j$, there are $\lambda_J(j)$ eligible left
runs and $\mu_J(j)$ eligible right runs. Once its run is chosen, its
position is forced by increasing order. Marking every left label by
$z$ gives the product in \eqref{eq:t97cut}; the lucky labels on the
left contribute $z^{|J|}$, and ordering the runs gives
$|J|!(k-|J|)!$. Conversely, the outcome recovers all these runs,
their order, and every assignment uniquely. Extracting the required
left length proves the formula, including empty left or right words.
\end{proof}

At $s=1$ only $J=\varnothing$ contributes, recovering the source's
Theorem~3.11:
\[
 c_1(n,I)=k!\prod_{j\notin I}(1+|I\cap[j-1]|).
\]
At $s=n+1$ only the all-left assignments contribute, recovering the
classical formula $k!\prod_{j\notin I}|I\cap[j-1]|$.
For $n=6$ and $I=\{1,3,5\}$, the full list for $s=1,\ldots,7$ is
\[
                 (144,72,82,82,82,78,36).
\]

\paragraph{Scope, prior methods, and verification.}
These formulas answer the two stated enumeration questions with their
original parameter ranges. In particular, general capacities are not
replaced by the capacity-one case. The first formula is based on an
auxiliary terminal-event enumeration and the source's established
outcome criterion. The second extends the increasing-run argument of
the source's Theorems~2.2 and~3.11. The precise formulas were not located
in a bounded primary-source search through 27 October 2026; this does
not certify worldwide priority. Neither terminal-order decomposition
nor the increasing-run principle is claimed to be historically new.
Accordingly, this track is not certified as meeting every strict
novelty condition in the request.

The standard-library program \path{t97_verify.py} checks
\eqref{eq:t97general} against direct ordered partitions in 728 cases
through $n=6$, covering every capacity composition and every adjacency
pattern. It independently runs the original parking procedure for 125
capacity vectors through $n=4$, checking the entire family of possible
lucky sets for each outcome. It also checks \eqref{eq:t97cut} in 1,792
fixed-set/spot cases through $n=7$. These finite checks supplement the
bijective all-parameter proofs above.

\clearpage
\section{T98: the exact shifted Rudin--Shapiro Hankel period}
\label{sec:t98}

The unchanged target is Conjecture~4.2 of Han,
\emph{Hankel continued fraction and its applications},
\emph{Advances in Mathematics} \textbf{303} (2016), 295--321
\cite{t98han2016}.
The conjecture concerns the binary sequence and determinants over
$\mathbb F_2$, not the signed Rudin--Shapiro sequence over the integers.
Let $u_n\in\mathbb F_2$ be the parity of the number of occurrences of
$11$ in the binary expansion of $n$. Set
\[
 f_k(x)=\sum_{n\ge0}u_{n+k}x^n,
 \qquad h_k(n)=\det(u_{k+i+j})_{0\le i,j<n},\qquad h_k(0)=1.
\]

\begin{theorem}[Han's original least-period formula]
For every $k\ge3$, let $m\ge1$ be determined by
$2^m<k\le2^{m+1}$. Then the sequence $(h_k(n))_{n\ge0}$ is
periodic from its first term, with least period
\[
                         R_k=9\cdot2^m.
\]
\end{theorem}

The auxiliary problem is to recognize unusually long gaps between normal
Pad\'e indices from the norm of an approximation error. In this problem,
a square-congruence at one finite place forces the norm to be a square;
odd pole orders at the only other possible finite pole then force that
norm to be exactly one. This identifies the long gaps with explicit
units. Their spacing proves \emph{minimality} of the Hankel period,
which does not follow from a continued-fraction period alone.

\subsection{The quadratic field and the shifted lattices}

Work throughout in characteristic two. Put $z=x^{-1}$, $y=z+1$, and
\[
 F_k(z)=z^{-1}f_k(z^{-1}),\qquad
 P_k(z)=\sum_{i=0}^{k-1}u_i z^{k-1-i}.
\]
Then $F_k=z^kF_0+P_k$. Han's quadratic equation for $f_0$ is
\[
 x^3+(1+x)^4f_0+(1+x)^5f_0^2=0.
\]
Consequently $g=yF_0$ satisfies
\[
                         g^2+g=y^{-3}.
\]
The function field $K=\mathbb F_2(z,g)$ has involution
$\overline g=g+1$. In particular,
\begin{equation}\label{eq:t98trace}
 \operatorname{Tr}(F_k)=\frac{z^k}{y},\qquad
 \operatorname{N}(F_0)=y^{-5}.
\end{equation}
The chosen expansion of $F_k$ at $z=\infty$ has strictly negative
degree, whereas its conjugate has degree $k-1$ for $k\ge2$.

There is only one possible finite pole of $F_0$: the point $Q$ above
$y=0$. In the variable $t=y^{-1}$ the curve is the smooth cubic
$g^2+g=t^3$. At its point at infinity,
\[
                v_Q(y)=2,\qquad v_Q(g)=-3,\qquad v_Q(F_0)=-5.
\]
Thus, if $A,B\in\mathbb F_2[z]$ and $A+BF_0$ has a finite pole,
that pole is at $Q$ and its order is one of $5,3,1$.
Indeed, a negative valuation must equal $2\operatorname{ord}_yB-5$;
the polynomial $A$ cannot cancel it.

Let $P_+$ and $P_-$ be the points $(t,g)=(0,0)$ and $(0,1)$.
They are the two places above $z=\infty$. The function
\[
 \varepsilon=\frac{g}{g+1}=1+y^3g=1+y^4F_0
\]
has norm one, trace $y^3$, and divisor $3P_+-3P_-$. Its degree in the
chosen Laurent expansion is $-3$.
Every norm-one function with no finite pole is a power of
$\varepsilon$. It also has no finite zero: its norm is one, and a
finite zero would require a finite pole of its conjugate. Thus its
divisor must be $j(P_+-P_-)$.
Since the curve is a smooth cubic of genus one, $P_+-P_-$ is not
principal: a function with one simple pole would give a degree-one map
to the projective line. Its order is exactly three because
$3(P_+-P_-)$ is the divisor just exhibited. Hence $3\mid j$;
division by the appropriate power of $\varepsilon$ leaves a constant,
and the only nonzero constant in $\mathbb F_2$ is one.

\subsection{Units that preserve a shifted approximation lattice}

Write $L_j(T)=\varepsilon^j+\varepsilon^{-j}$ when
$T=\varepsilon+\varepsilon^{-1}$. These polynomials satisfy
$L_0=0$, $L_1=T$, and $L_{j+1}=TL_j+L_{j-1}$. For $j>0$,
\begin{equation}\label{eq:t98orders}
 \operatorname{ord}_z L_j(y^3)=
 \begin{cases}
  0,&3\nmid j,\\
  2^{\nu_2(j)+1},&3\mid j,
 \end{cases}
 \qquad
 \operatorname{ord}_yL_j(y^3)=3\cdot2^{\nu_2(j)}.
\end{equation}
Here $\nu_2$ is the exponent of two in $j$. For the first formula,
$L_j(1)=0$ precisely when $3\mid j$, and
$L_3(T)=T(T+1)^2$, while $y^3+1=z(z^2+z+1)$.
For odd $a$, $L_a(X)=X+O(X^3)$; use composition
$L_a(L_b(T))=L_{ab}(T)$ and the Frobenius identity
$L_{2a}=L_a^2$. These facts prove both formulas.

Fix $k\ge3$ and put $J=3\cdot2^m$ with $2^m<k\le2^{m+1}$.
For $U=\varepsilon^J$, define
\[
 b=\frac{yL_J(y^3)}{z^k},\qquad a=U+bF_k.
\]
By \eqref{eq:t98orders}, $b$ is a polynomial and is divisible by
$y^7$. The function $a$ is invariant under the involution, so belongs
to $\mathbb F_2(z)$; it has no finite pole and is therefore a polynomial.
Moreover $b\operatorname{Tr}(F_k)$ and
$b\operatorname{N}(F_k)$ are polynomials. The latter norm has denominator
dividing $y^5$, and $b$ is divisible by $y^7$.
Multiplication by $U$ on
\[
             \mathcal M_k=\mathbb F_2[z]+\mathbb F_2[z]F_k
\]
therefore has the polynomial matrix
\begin{equation}\label{eq:t98matrix}
 \begin{pmatrix}
 a&b\operatorname{N}(F_k)\\
 b&a+b\operatorname{Tr}(F_k)
 \end{pmatrix}.
\end{equation}
Its determinant is $\operatorname{N}(U)=1$, so it is unimodular and
preserves coprime pairs of coefficients in both directions.
The same argument applies to every power $U^\ell$.

Conversely, if $\varepsilon^j=p+qF_k$ with polynomial $p,q$, taking
traces forces
\[
 q=\frac{yL_j(y^3)}{z^k}.
\]
Thus \eqref{eq:t98orders} implies $J\mid j$.
This also proves that $J$ is the least positive exponent of a unit in
$\mathcal M_k$.

\subsection{A period for all normal indices}

We use the standard formal continued-fraction correspondence: the normal
indices $n$, meaning $h_k(n)\ne0$, are exactly the denominator degrees of
the convergents to $F_k$, including the initial convergent $0/1$ of
degree zero. For a reduced proper convergent $p/q$, with $q$ monic and
$\deg q=n$, write
\[
                    \deg(qF_k+p)=-(n+d),\qquad d\ge1.
\]
Then $d$ is the gap to the next normal index. Equivalently, a reduced
fraction is a convergent precisely when
$\deg(qF_k+p)<-\deg q$. These are the usual formal Pad\'e/continued-
fraction facts (the strict approximation criterion is Proposition~1 of
van der Poorten~\cite{t98vdp2001}); Han's Hankel-fraction Theorem~2.1~\cite{t98han2016}
also identifies the
normal indices as cumulative positive gaps. The nonzero determinants
are all one because the coefficient field is $\mathbb F_2$.

Let $\alpha=p+qF_k$ represent such a convergent. Its conjugate has
degree $n+k-1$, by \eqref{eq:t98trace}. Multiplication by $U$ increases
the vanishing order of $\alpha$ at $P_+$ by
$R=3J=9\cdot2^m$, and increases the degree of its conjugate by $R$.
By the unimodular matrix \eqref{eq:t98matrix}, the resulting
$p'+q'F_k$ has coprime polynomial coefficients. Taking its trace shows
$\deg q'=n+R$, while its degree is $-(n+R+d)$. It is therefore a
convergent of degree $n+R$ with the same following gap $d$.

Conversely, multiply a convergent of denominator degree $n+R$ by
$U^{-1}$, for $n\ge0$. The same degree calculation gives a reduced
convergent of degree $n$ with its gap unchanged. Hence
\[
                   h_k(n+R)=h_k(n)\qquad(n\ge0).
\]
This proves a period from the first term, including cases where some
Hankel determinants vanish. It does not yet prove least period.

\subsection{The long-gap norm criterion and least period}

The norm argument solves the following auxiliary problem in an entire
family of quadratic function fields.

\begin{lemma}[Odd-pole norm rigidity]\label{lem:t98rigidity}
Let $s\ge1$, $y=z+1$, and $\phi=g/y$, where
$g^2+g=y^{-(2s-1)}$ over $\mathbb F_2(z)$. Let $A,B$ be polynomials,
not both zero, with $z^k\mid B$ and $k\ge2s+2$. If
$\deg\operatorname{N}(A+B\phi)\le0$, then
\[
                         \operatorname{N}(A+B\phi)=1.
\]
\end{lemma}
\begin{proof}
Write
\[
 N=A^2+\frac{AB}{y}+\frac{B^2}{y^{2s+1}}
    =\frac{V}{y^{2s+1}},\qquad\deg V\le2s+1.
\]
At $z=0$, $N\equiv A^2\pmod{z^k}$, so $N'$ vanishes to order
at least $k-1$. But
\[
 N'=\frac{yV'+V}{y^{2s+2}},\qquad\deg(yV'+V)\le2s.
\]
As $k-1>2s$, this derivative is zero and $N$ is a rational square.
The odd-order pole of $y^{-(2s-1)}$ is totally ramified in the
Artin--Schreier extension. Its unique point $Q$ has
$v_Q(y)=2$, $v_Q(g)=-(2s-1)$, and $v_Q(\phi)=-(2s+1)$.
If $A+B\phi$ had a finite pole, its valuation would therefore be a
negative odd number $2\operatorname{ord}_yB-(2s+1)$.
The norm has the same valuation at $y=0$, since the residue degree
is one, contradicting its being a square. There are no other possible
finite poles. Hence $N$ is a nonzero polynomial of degree at most zero,
and so equals one.
\end{proof}

\begin{lemma}[A long gap forces a norm-one error]
Let $k\ge6$, and let $p/q$ be a reduced convergent of denominator
degree $n$, with following normal-index gap $d\ge k-1$. Then
\[
                       p+qF_k=\varepsilon^{\ell J}
\]
for some positive integer $\ell$. Conversely, each of these positive
unit powers is such an approximation error. Its gap is exactly $k-1$,
and the next normal index is $\ell R$.
\end{lemma}

\begin{proof}
Put $\alpha=p+qF_k=A+BF_0$, where
$A=p+qP_k$ and $B=qz^k$ are polynomials. Its norm is
\[
 N=\operatorname{N}(\alpha)
   =A^2+\frac{AB}{y}+\frac{B^2}{y^5}
   =\frac{V(z)}{y^5}.
\]
Its degree is $(n+k-1)-(n+d)=k-1-d\le0$.
Apply Lemma~\ref{lem:t98rigidity} with $s=2$ to obtain $N=1$.
The proof of that lemma also shows that $\alpha$ has no finite pole,
so the unit
classification gives $\alpha=\varepsilon^j$; its positive vanishing
order at $P_+$ gives $j>0$. Polynomial membership gives $J\mid j$.

For the converse, the polynomial matrix for $\varepsilon^{\ell J}$
shows that its coefficients $p,q$ in $\mathcal M_k$ are coprime.
Taking traces gives
\[
                \deg q=\ell R-k+1,
                \qquad\deg\varepsilon^{\ell J}=-\ell R.
\]
Thus it is a reduced convergent, with gap $k-1$ and next denominator
degree $\ell R$, as asserted.
\end{proof}

For $k\ge6$, the endpoints of all gaps at least $k-1$ are now exactly
$R,2R,3R,\ldots$. Every period of the binary Hankel sequence must
preserve this set of gap endpoints sufficiently far from its initial
index, so it must be divisible by $R$. Since $R$ is already a period,
it is the least period.

\subsection{The three small shifts}

For completeness, $k=3,4,5$ admit short exact quadratic-transform
certificates. Encode a polynomial $\sum a_i x^i\in\mathbb F_2[x]$ by
the integer $\sum a_i2^i$. Here the initial power-series tail is the
original $f_k(x)$, not the Cauchy series $F_k(z)$. For a current tail
$W(x)$ satisfying $A+BW+CW^2=0$ with
$B=1+x^4$, let $d=\operatorname{ord}_x A$ and let $D$ be the
degree-$(d+1)$ truncation of $x^d/W$. Substitution
$W=x^d/(D+x^{d+2}G)$ gives the exact update
\[
 (A,B,C)\longmapsto
 \left(\frac{AD^2+x^dBD+x^{2d}C}{x^{2d+2}},\ B,\ x^2A\right).
\]
The following lists of pairs $(d+1,D)$ return exactly to their initial
triple, so their cumulative first coordinates specify all normal indices.
\begin{center}\small
\begin{tabular}{ccl}
\toprule
$k$ & Initial $(A,C)$ & Cycle pairs $(d+1,D)$\\
\midrule
3 & $(1,408)$ & $(1,1),(2,1),(1,1),(1,3),(1,3),(1,1)$\\
  & & $(1,3),(2,1),(1,3),(3,1),(2,7),(2,7)$\\[2pt]
4 & $(196,816)$ & $(3,1),(2,3),(1,1),(1,1),(3,1)$\\
  & & $(1,3),(2,1),(1,1),(1,3),(3,15)$\\[2pt]
5 & $(98,1632)$ & $(2,1),(3,11),(2,5),(2,1),(2,3),(1,1)$\\
  & & $(2,5),(1,3),(2,7),(4,1),(1,3),(2,5)$\\
  & & $(4,1),(1,3),(1,1),(1,3),(1,1),(4,31)$\\
\bottomrule
\end{tabular}
\end{center}
Their binary Hankel period words, beginning at $n=0$, are respectively
\[
 \begin{array}{c|l}
 3&110111111011001010\\
 4&100101110011011100\\
 5&101001010101101101000110100011111000.
 \end{array}
\]
Checking the proper divisors of their lengths shows that their least
periods are $18,18,36$. This completes the all-$k$ proof.

\paragraph{Provenance, novelty, and verification.}
The original conjecture, its coefficient field, and its least-period
requirement are retained. The proof above supplies both the period and
its minimality. The new proposed argument is the explicit shifted-lattice
unit computation combined with the long-gap norm criterion. Quadratic
function fields, formal continued fractions, units, and genus-one
divisors are established tools; a historically unprecedented foundational
principle is not certified here. In particular, the broad connection
between norm-one units, large partial quotients, and continued-fraction
quasi-periods is classical; see van der Poorten~\cite[Sections 2--3]{t98vdp2001}.
The characteristic-two odd-pole norm criterion and its application to the
precise shifted lattice are the additional arguments used here.
A bounded search located no prior
resolution of this precise conjecture, but that is not a guarantee of
worldwide priority. The accompanying verifier checks the finite cycles,
unit matrices, trace orders, and independent small Hankel determinants.
The infinite conclusion follows from the proof, not from these finite
checks.

\clearpage
\section{T99: exact rank-two average mixing on 128 and 256 vertices}
\label{sec:t99}

The unchanged external target is Open Problem~4.1 in the final published
version of Coutinho--Guo \cite{t99cg2024}: whether there are only finitely
many graphs whose average mixing matrix has rank two. The numbering here
is that of the published article, not its first arXiv version.
\emph{This investigation does not answer that finiteness question.}
It constructs two explicit connected simple graphs of orders 128 and 256
and proves that their average mixing matrices have rank exactly two.
The missing step is an unbounded construction, or a finiteness theorem.

For a real symmetric adjacency matrix $A$, write
$A=\sum_\theta\theta E_\theta$ for its spectral decomposition. Its average
mixing matrix is
\[
 \widehat M=\lim_{R\to\infty}\frac1R\int_0^R
       e^{itA}\circ e^{-itA}\,dt
       =\sum_\theta E_\theta\circ E_\theta,
\]
where $\circ$ denotes entrywise multiplication. The final equality follows
by integrating each exponential $e^{it(\theta-\tau)}$.

\subsection{An auxiliary construction problem with an exact answer}

Let $G=\mathbb F_2^d$, $N=2^d$, and choose
$S,T\subseteq G\setminus\{0\}$ and $B\subseteq G$.
Define a graph $X(S,B,T)$ on $G\times\{0,1\}$ by the edges
\[
 (g,0)\sim(g+s,0),\quad
 (g,1)\sim(g+t,1),\quad
 (g,0)\sim(g+b,1),
 \qquad s\in S,\ t\in T,\ b\in B.
\]
Every group element is its own inverse, so this is an undirected simple
graph. For $x\in G$, put $\chi_x(g)=(-1)^{x\cdot g}$ and
\[
 a_x=\sum_{s\in S}\chi_x(s),\qquad
 b_x=\sum_{b\in B}\chi_x(b),\qquad
 c_x=\sum_{t\in T}\chi_x(t),
\]
\[
 t_x=a_x+c_x,\qquad u_x=a_x-c_x,\qquad
 D_x=u_x^2+4b_x^2.
\]
The auxiliary question is to give an exact finite certificate, using only
these integer sums, for simple adjacency spectrum and rank-two mixing.

\begin{proposition}[Parity certificate and exact mixing formula]
\label{prop:t99certificate}
Suppose $|S|$ is even, $|T|$ and $|B|$ are odd, and the $N$ ordered pairs
$(t_x,D_x)$ are distinct. Then $A(X(S,B,T))$ has $2N$ distinct eigenvalues,
and
\begin{equation}\label{eq:t99mixing}
 \widehat M=\frac{J_{2N}}{2N}+\frac{\kappa}{2N}ss^{\mathsf T},
 \qquad s=(\mathbf1_N,-\mathbf1_N)^{\mathsf T},\qquad
 \kappa=\frac1N\sum_{x\in G}\frac{u_x^2}{D_x}.
\end{equation}
Here $0<\kappa<1$, and $\operatorname{rank}\widehat M=2$.
If $S$ spans $G$, the graph is connected.
\end{proposition}

\begin{proof}
For each $x$, the two-dimensional subspace consisting of $\chi_x$ on
either fiber is invariant under adjacency. In its normalized character
basis adjacency acts as
\[
 H_x=\begin{pmatrix}a_x&b_x\\b_x&c_x\end{pmatrix},\qquad
 p_x(z)=\det(zI-H_x)=z^2-t_xz+a_xc_x-b_x^2.
\]
All $N$ such subspaces are mutually orthogonal and together exhaust the
vertex space. This is the standard character decomposition of a
semi-Cayley graph \cite{t99gl2010}.

Since a sum of signs has the parity of its number of summands, $u_x$ and
$b_x$ are odd for every $x$. Thus
$D_x\equiv1+4\equiv5\pmod8$. In particular $D_x>0$ is not a square,
and $p_x$ is irreducible over $\mathbb Q$. Distinct pairs $(t_x,D_x)$
give distinct monic polynomials $p_x$. Two distinct monic irreducible
quadratics over $\mathbb Q$ cannot have a common root. Their roots
$(t_x\pm\sqrt{D_x})/2$ therefore give $2N$ distinct eigenvalues.

Put
\[
 K_x=2H_x-t_xI=\begin{pmatrix}u_x&2b_x\\2b_x&-u_x\end{pmatrix}.
\]
Then $K_x^2=D_xI$, so the two spectral projectors in this block are
$P_{x,\pm}=(I\pm K_x/\sqrt{D_x})/2$.
For fiber-first ordering of the vertices the full projector is
$P_{x,\pm}\otimes(\chi_x\chi_x^{\mathsf T}/N)$.
Since every entry of $\chi_x$ has square one,
\[
 \widehat M=\frac1{N^2}\sum_x
 (P_{x,+}\circ P_{x,+}+P_{x,-}\circ P_{x,-})\otimes J_N.
\]
The diagonal entries of the parenthesized $2\times2$ matrix are
$(1+u_x^2/D_x)/2$; its off-diagonal entries are
$2b_x^2/D_x=(1-u_x^2/D_x)/2$.
Summing proves \eqref{eq:t99mixing}. The nonzero orthogonal vectors
$\mathbf1_{2N}$ and $s$ have eigenvalues $1$ and $\kappa$, respectively;
all vectors orthogonal to both have eigenvalue zero. Oddness makes both
$u_x$ and $b_x$ nonzero, whence $0<u_x^2/D_x<1$ and $0<\kappa<1$.
This proves the rank assertion.

If $S$ spans $G$, the graph induced on fiber zero is connected.
Oddness makes $B$ nonempty, and any fixed $b\in B$ joins every vertex
of fiber one to a vertex of fiber zero. The full graph is connected.
\end{proof}

\subsection{Two explicit graphs and their finite proof certificates}

Identify $G$ with the integers $0,\ldots,2^d-1$ using binary coordinates,
so addition is bitwise XOR. For a nonnegative integer $h<2^N$, define
\[
 \mathcal S(h)=\{g\in\{0,\ldots,N-1\}:
                \lfloor h/2^g\rfloor\equiv1\pmod2\}.
\]
The following table gives the hexadecimal integers $h_S,h_B,h_T$;
the sets are exactly $S=\mathcal S(h_S)$, $B=\mathcal S(h_B)$,
$T=\mathcal S(h_T)$. Thus the specification is complete and does not
depend on a random seed or on reproducing the search.

\begin{center}
\begin{tabular}{ccl}
\toprule
$d$ & Set & Hexadecimal membership mask\\
\midrule
6 & $S$ & \texttt{1d39f52fe7384174}\\
6 & $B$ & \texttt{37ec8fd2075e6984}\\
6 & $T$ & \texttt{29e5fdffe09a47da}\\
\midrule
7 & $S$ & \texttt{623a7e9d385f8df96b33f5d522a56aa8}\\
7 & $B$ & \texttt{db196b1a5d26bcf475c59b8a07829823}\\
7 & $T$ & \texttt{c8975afdc8d54101d50561d3475eab7e}\\
\bottomrule
\end{tabular}
\end{center}

\begin{proposition}[Two finite examples]
The displayed masks define connected simple graphs on 128 and 256
vertices, respectively, with simple adjacency spectrum and average
mixing rank exactly two.
\end{proposition}

\begin{proof}[Exact finite verification]
The accompanying dependency-free program \texttt{t99\_verify.py}
decodes the six integers above. For $d=6$, it gives
$(|S|,|B|,|T|)=(34,33,39)$; for $d=7$, it gives $(70,63,65)$.
In both cases $0\notin S\cup T$ and Gaussian elimination over
$\mathbb F_2$ gives $\dim\langle S\rangle=d$.
For each of the 64 or 128 characters the program computes $a_x,b_x,c_x$
in two ways: direct signed sums and an integer Walsh transform. These
agree exactly. It then checks that $D_x\equiv5\pmod8$ and that every
pair $(t_x,D_x)$ is distinct. These are finite integer assertions, with
no floating-point eigenvalue calculation. The program also independently
constructs all edges and checks the graph's simplicity and connectivity.
All assertions pass. Proposition~\ref{prop:t99certificate} supplies the
conclusion. The file \texttt{t99\_verify.json} records the expanded sets,
every quadratic polynomial's trace, determinant and discriminant, and
the exact rational value of $\kappa$ for independent inspection.
\end{proof}

The two fiber degrees are respectively $(67,72)$ and $(133,128)$.
The construction and verifier certify these two graphs individually;
neither numerical rank estimation nor an extrapolation from small
orders is used in the proof.

\subsection{The obstruction to the attempted infinite construction}

One attempted recursion transformed a column vector of three Walsh
coefficients by either
\[
 M_+=\begin{pmatrix}0&1&0\\1&0&1\\1&1&0\end{pmatrix},\qquad
 M_-=\begin{pmatrix}0&1&0\\-1&0&1\\1&-1&0\end{pmatrix}.
\]
It cannot yield a collision-free family at every depth. Applying the
operations in either of the two orders
$-+++--$ and $+---++$ gives exactly the same matrix product,
\[
 \begin{pmatrix}1&0&0\\-4&1&0\\0&0&1\end{pmatrix}.
\]
Consequently two distinct branches give identical coefficient triples
after six steps, for every initial vector. This rejects that specific
recursion; it is not an obstruction to all possible constructions.

The spectral decomposition used above is classical, and the rank bound
also has a general explanation: with simple spectrum, a graph
automorphism sends each real normalized eigenvector to its positive or
negative. Squared eigenvector entries are therefore constant on each
automorphism orbit. If a subgroup has two orbits, the squared-eigenvector
matrix has rank at most two, and so does its Gram matrix $\widehat M$.
The average-mixing framework is also established in
Coutinho--Godsil--Guo--Zhan \cite{t99cggz2018}.

Accordingly, the contribution here is the explicit finite construction,
its parity certificate and exact mixing formula, together with a
checkable failure of the proposed recurrence. A bounded literature
search did not locate the identical examples or an answer to the
published finiteness question. That search does not establish worldwide
priority. No historically new foundational principle, infinite family,
or complete solution of Open Problem~4.1 is asserted.

\clearpage
\section{T100: the nonvanishing gap in Sun's quadratic-residue determinant}
\label{sec:t100}

The unchanged target is Conjecture~4.10(ii) of Sun
\cite{t100sun2024}. For a prime $p\equiv3\pmod4$, put $n=(p-1)/2$
and write $\chi$ for the Legendre symbol, with $\chi(0)=0$. The
conjecture is
\[
 \chi\!\left(2\det A_p\right)=1,
 \qquad
 A_p=\bigl[\chi(j+k)+\chi(j^2+k^2)\bigr]_{0\le j,k\le n}.
\]
The factor $2$ is outside the determinant. All indices and the prime
range in this statement are retained.

Jiang--Sun \cite[Theorem~1.1 and Remark~1.1]{t100jiangs2026}
already prove that the displayed Legendre symbol is $0$ or $1$, and
explicitly state that they cannot prove $p\nmid\det A_p$. Thus merely
reproducing a square factorization would not resolve the conjecture.
The present investigation also does \emph{not} prove this nonvanishing.
We give the precise surviving obstruction and an all-prime reason why
a proposed extension to arbitrary half-systems cannot supply it.

\subsection{The exact existing reduction}
All calculations in this subsection take place over $\mathbb F_p$.
Let
\[
 P(X)=\prod_{k=0}^{n}(X-k),\qquad
 r_m(X)=X^m\bmod P(X),\qquad q=(n+1)/2.
\]
Define the $q\times q$ matrix
\[
 C_{ab}=
 \mathbf1_{a+b=q-1}\binom n{2b+1}
 +\binom na[X^{2b+1}]r_{2n-2a}(X),
 \qquad 0\le a,b<q.
\]
The parity-block argument in Jiang--Sun gives
\begin{equation}\label{eq:t100compression}
 \det A_p=(-1)^q(\det V)^2(\det C)^2,
 \qquad V=(j^u)_{0\le j,u\le n}.
\end{equation}
For clarity, Euler's criterion replaces the entries by
$(j+k)^n+(j^2+k^2)^n$. Reducing each variable modulo $P$ produces
a bivariate polynomial of degree at most $n$ in each variable. Its
coefficient matrix has zero odd--odd block: the first summand has odd
total degree, and in each term of the second summand at least one
exponent is even and at most $n-1$. Ordering even exponents before odd
ones gives a symmetric block matrix of the form
\[
 \begin{pmatrix} B&C\\ C^T&0\end{pmatrix}.
\]
Multiplication by $V$ and $V^T$ gives the evaluated matrix, proving
\eqref{eq:t100compression}. This is the published reduction, credited
here rather than claimed as new.

Since the elements $0,1,\ldots,n$ are distinct, $\det V\ne0$.
Also $\chi(2)(-1)^q=1$. Consequently the original missing assertion is
\emph{exactly}
\[
                         \det C\ne0\quad\hbox{in }\mathbb F_p.
\]
This is a reduction of the original target, not a proved replacement
for it.

\subsection{An auxiliary question and an obstruction to a general principle}

A half-system is a set $H=\{0,h_1,\ldots,h_n\}\subset\mathbb F_p$
containing exactly one element of every nonzero pair $\{a,-a\}$.
For such a set define
\[
 P_H(X)=\prod_{h\in H}(X-h)=E_H(X^2)+XO_H(X^2),
 \qquad
 A_H=\bigl[\chi(h+k)+\chi(h^2+k^2)\bigr]_{h,k\in H}.
\]
The proposed auxiliary question was whether the common norm identity
for $E_H,O_H$, together with the half-system property, forces
$\det A_H\ne0$. It has the following complete negative answer.

\begin{proposition}[A universal obstruction for half-systems]
For every prime $p\equiv3\pmod4$ and every half-system $H$,
\begin{equation}\label{eq:t100pell}
 E_H(Z)^2-ZO_H(Z)^2=Z^{n+1}-Z.
\end{equation}
Nevertheless, for every such prime there is a half-system for which
$A_H$ has an identically zero row and hence zero determinant even over
the integers. One may take $H$ to be zero together with the quadratic
nonresidues modulo $p$.
\end{proposition}

\begin{proof}
The degree $n+1$ is even. The roots of $P_H(X)P_H(-X)$ are every
nonzero field element once and zero twice, and its leading coefficient
is one. Hence
\[
 P_H(X)P_H(-X)=X\prod_{a\in\mathbb F_p}(X-a)
             =X^{p+1}-X^2.
\]
Taking even and odd parts gives \eqref{eq:t100pell}.

Since $\chi(-1)=-1$, exactly one member of each pair $\{a,-a\}$ is
a quadratic nonresidue. The proposed $H$ is therefore a half-system.
For each nonzero $k\in H$, the entry in row zero is
\[
                     \chi(k)+\chi(k^2)=-1+1=0.
\]
The entry at $(0,0)$ is zero as well. In this example
$P_H(X)=X(X^n+1)$, so $E_H(Z)=Z^q$ and $O_H(Z)=1$ explicitly satisfy
the same norm identity. Thus that identity does not prevent singularity.
\end{proof}

This is a counterexample to the auxiliary extension, \emph{not} a
counterexample to Sun's original conjecture. Sun specifies the
consecutive half-system $H_0=\{0,1,\ldots,n\}$, and replacing it by
the nonresidues would change the problem.

The consecutive set supplies additional identities that survive the
obstruction above:
\begin{equation}\label{eq:t100interval}
 P(-X-1/2)=P(X),\qquad
 (X+1/2)P(X+1)=(X+1)P(X).
\end{equation}
Indeed, $n=-1/2$ in $\mathbb F_p$, reflection $k\mapsto n-k$ permutes
the roots, and the second identity follows by shifting the product
defining $P$. The first identity also says that $P$ is a polynomial in
$(X+1/4)^2$. These interval identities distinguish the required set
from an arbitrary half-system, but no triangular factorization or other
nonvanishing proof for $C$ was obtained from them.

\paragraph{Verification and success status.}
The exact verifier \texttt{t100\_verify.py} checks the original matrix,
the coefficient matrix, \eqref{eq:t100compression}, both interval
identities, the norm identity, and the nonresidue obstruction for all
primes $p<200$ with $p\equiv3\pmod4$. All original determinants in
that finite range are nonzero modulo $p$. This is finite evidence and
does not establish the all-prime assertion. T100 remains an unsuccessful
attempt at the unchanged named problem. The auxiliary negative result
and the existing compression formula are not counted as a new complete
solution or as a historically unprecedented method.
\clearpage
\section*{Original-target status and reproducibility}
\addcontentsline{toc}{section}{Original-target status and reproducibility}

The scope of each conclusion is recorded below. ``Complete'' refers to
the written proof of the exact source question, not to a certification
of historical priority or to external acceptance.

\begin{center}\small
\begin{tabular}{@{}>{\raggedright\arraybackslash}p{0.08\textwidth}>{\raggedright\arraybackslash}p{0.35\textwidth}>{\raggedright\arraybackslash}p{0.50\textwidth}@{}}
\toprule
Track & Pre-existing target & Conclusion of this manuscript\\
\midrule
T96 & Ulas, Question 2.10 & Partial: residue-one lower bound with exponent
$\log_2 3<2$; the requested limit is not proved.\\[5pt]
T97 & Ferreri--Harris--Martinez--Swartz, Section 6, first two questions &
Complete coefficient formulas over the original parameter ranges.
These questions have no individual conjecture numbers.\\[5pt]
T98 & Han, Conjecture 4.2 & Complete all-shift proof of the exact least
period, using the odd-pole norm lemma and explicit units.\\[5pt]
T99 & Coutinho--Guo, Open Problem 4.1 & Exact finite rank-two examples;
the finite-versus-infinite question remains unresolved.\\[5pt]
T100 & Sun, Conjecture 4.10(ii) & An auxiliary obstruction is proved;
nonvanishing for the prescribed consecutive set remains unresolved.\\
\bottomrule
\end{tabular}
\end{center}

The main distinction in T98 is the additional characteristic-two
argument in Lemma~\ref{lem:t98rigidity}, followed by the precise
classification of long gaps for the shifted lattices. Its use of
classical function-field and continued-fraction theory is explicit.
T97 uses classical combinatorial decompositions. Neither track is
described as a historically unprecedented foundational principle.

\paragraph{Literature scope.}
The source and priority review used exact paper identifiers, author and
conjecture numbers, distinctive formulas, and related primary papers,
through 27 October 2026. In particular, the shifted binary
Rudin--Shapiro problem was separated from signed Rudin--Shapiro
nonsingularity, general automaticity of Hankel arrays, and other
quadratic-series periodicity results. The parking counts were separated
from counts of preference lists. The average-mixing target was kept as
the ordinary Cesaro average, and the nonzero requirement in Sun's
conjecture was separated from the published square-or-zero result.
No identical resolution of the two T97 questions or Han's precise T98
conjecture was located in that bounded review. This is negative search
evidence, not a proof that no such literature exists. An unretrieved
potentially related thesis lead in the parking search was not treated
as cleared. The verification archive includes the detailed source audit.

\paragraph{Reproduction.}
The archive contains this standalone \LaTeX{} source, the compiled PDF,
five exact-arithmetic Python verifiers, their JSON results, a single
runner, and a file-hash manifest. The programs use Python's standard
library and do not require floating-point eigenvalue tests, symbolic
algebra packages, or network access. Run
\begin{center}
\texttt{python3 run\_all.py}
\end{center}
from the extracted archive, without Python's \texttt{-O} option.
The runner checks file hashes before execution and reports each
program's exit status. Each individual verifier can also be run directly.
The \LaTeX{} file compiles with pdf\LaTeX{} using standard AMS,
Latin Modern, geometry, microtype, and hyperref packages.

Finite computations check the displayed certificates and independent
instances. They do not prove the unbounded parameter assertions by
themselves. The written bijections, degree arguments, function-field
proof, and explicit counterexample construction supply the infinite
conclusions stated in the respective sections. Internal independent
mathematical reviews found no unresolved gap in those claimed proofs;
this is not external peer review or proof-assistant certification.

\paragraph{Strict success condition.}
The demand that all five tracks satisfy every historical and mathematical
condition is not met. T98 supplies a complete proof of an unchanged,
pre-existing named conjecture as stated in its original source; T97
supplies two complete formulas for explicit source questions. The
remaining three original targets are not declared solved. Worldwide
priority and a historically unprecedented principle remain unverified,
including for the complete proofs.

\begin{thebibliography}{99}
\raggedright
\bibitem{t96:ulas}
M.~Ulas, \emph{Strong arithmetic property of certain Stern polynomials},
Publicationes Mathematicae Debrecen \textbf{96} (2020), nos.~3--4;
arXiv:1909.10844v1, Corollary~2.7 and Question~2.10.
\url{https://arxiv.org/html/1909.10844v1}.

\bibitem{t96:spiegelhofer}
L.~Spiegelhofer, \emph{A digit reversal property for Stern polynomials},
Integers \textbf{17} (2017), Article A53; arXiv:1610.00108v2.
\url{https://arxiv.org/abs/1610.00108v2}.
\bibitem{t97fhms}
M.~Ferreri, P.~E.~Harris, L.~Martinez, and E.~Swartz,
\emph{Enumerating Vector Parking Functions and their Outcomes Based on
Specified Lucky Cars}, arXiv:2508.13917v2, 10 September 2025,
Theorems~2.2, 3.7, 3.11, Definition~3.14, and Section~6.
\url{https://arxiv.org/html/2508.13917v2}.
\bibitem{t98han2016}
G.-N.~Han, \emph{Hankel continued fraction and its applications},
Advances in Mathematics \textbf{303} (2016), 295--321,
Theorem~2.1 and Conjecture~4.2.
Author's full text, dated 22 August 2016:
\url{https://irma.math.unistra.fr/~guoniu/papers/p94hfrac.pdf}.

\bibitem{t98vdp2001}
A.~J.~van der Poorten, \emph{Non-periodic continued fractions in
hyperelliptic function fields}, Bulletin of the Australian Mathematical
Society \textbf{64} (2001), 331--343, Proposition~1 and Sections~2--3.
DOI: \href{https://doi.org/10.1017/S000497270003999X}{10.1017/S000497270003999X}.
\bibitem{t99cg2024}
G.~Coutinho and K.~Guo, \emph{Selected open problems in continuous-time
quantum walks}, Special Matrices \textbf{12} (2024), 20240025,
Open Problem~4.1 in the final published version.
\url{https://doi.org/10.1515/spma-2024-0025}.

\bibitem{t99gl2010}
X.~Gao and Y.~Luo, \emph{The spectrum of semi-Cayley graphs over abelian
groups}, Linear Algebra and its Applications \textbf{432} (2010),
2974--2983. \url{https://doi.org/10.1016/j.laa.2009.12.040}.

\bibitem{t99cggz2018}
G.~Coutinho, C.~Godsil, K.~Guo, and H.~Zhan,
\emph{A new perspective on the average mixing matrix},
Electronic Journal of Combinatorics \textbf{25}(4) (2018), Paper~4.14.
\url{https://arxiv.org/abs/1709.03591}.
\bibitem{t100sun2024}
Z.-W.~Sun, \emph{Problems and results on determinants involving Legendre
symbols}, arXiv:2405.03626v8, 28 July 2024, Conjecture~4.10(ii).
\url{https://arxiv.org/html/2405.03626v8}.

\bibitem{t100jiangs2026}
B.~Jiang and Z.-W.~Sun, \emph{Some new results on determinants and
permanents}, arXiv:2606.03970v1, 2 June 2026, Theorem~1.1,
Remark~1.1, and Section~2.
\url{https://arxiv.org/html/2606.03970v1}.

\end{thebibliography}
\end{document}
